% ==================== EXERCICE 2 : FONCTIONS ====================

\section{Fonctions \hfill \textit{6 points}}

Pour les fonctions suivantes, donner le \textbf{prototype}, la \textbf{définition} et un \textbf{exemple d'appel} :
    \UPSTIquestion Fonction \texttt{pair()} qui reçoit un entier et qui retourne 1 si cet entier est pair, 0 dans le cas contraire.
\espaceReponse{5cm}
    \UPSTIquestion Fonction \texttt{arrondi()} qui reçoit un réel positif et qui retourne l'entier le plus proche.
\espaceReponse{5cm}
    \UPSTIquestion Fonction \texttt{gereCompteur()} qui reçoit un entier (le compteur) et un booléen, et qui retourne le compteur modifié. Le compteur est incrémenté (+1) si le booléen reçu vaut 1, décrémenté (-1) si le booléen reçu vaut 0.

\espaceReponse{5cm}
